\subsection{Stability}

\begin{definition} \label{def:diagonal-dominance} \hlt{Diagonal Dominance}\\
Let $A=[a_{ij}]\in\operatorname{Mat}(n,n,\mathbb{C})$. Matrix $A$ is \hlt{diagonally dominant} if $\abs{a_{kk}}\geq R_k^{\mathrm{off}}(A)$ for every $k$, and \hlt{strictly diagonally dominant} if $\abs{a_{kk}}>R_k^{\mathrm{off}}(A)$ for every $k$.
\end{definition}

\begin{corollary} \label{cor:strict-diagonal-dominance-invertible} \hlt{Strict Diagonal Dominance Implies Invertibility}\\
Let $A\in\operatorname{Mat}(n,n,\mathbb{C})$.
\begin{enumerate}[label=(\roman*)]
\setlength{\itemsep}{0pt}
\item If $A$ is strictly diagonally dominant, then $A$ is invertible.
\item If $A$ has real spectrum and nonnegative diagonal entries, then diagonal dominance gives $\operatorname{spec}(A)\subseteq[0,\infty)$; strict diagonal dominance gives $\operatorname{spec}(A)\subseteq(0,\infty)$.
\end{enumerate}
\begin{proof}
\begin{enumerate}[label=(\roman*)]
\setlength{\itemsep}{0pt}
\item Strict dominance gives $0\notin \operatorname{Ger}_k(A)$ for every $k$, since $\abs{a_{kk}}>R_k^{\mathrm{off}}(A)$.\\
Thus $0\notin\operatorname{spec}(A)$ by Gersgorin, so $A$ is invertible by Corollary \ref{cor:matrix-invertible-iff-zero-not-spectrum}.
\item If $a_{kk}\geq0$ and $a_{kk}\geq R_k^{\mathrm{off}}(A)$, then each $\operatorname{Ger}_k(A)$ lies in the closed right half-plane. In the strict case, each disk lies in the open right half-plane. Gersgorin and the real-spectrum hypothesis give the claim.
\end{enumerate}
\hlt{Structural Commentary}: Gersgorin converts diagonal dominance into disk placement.\\
\hlt{Key Insight}: The disks control whether eigenvalues can reach $0$ or the left half-plane.
\end{proof}
\end{corollary}

\begin{corollary} \label{cor:diagonal-dominant-hermitian-positive-semidefinite} \hlt{Diagonal Dominance and Positive Semidefiniteness}\\
Let $A=[a_{ij}]\in\operatorname{Mat}(n,n,\mathbb{C})$ be Hermitian with nonnegative diagonal entries.
\begin{enumerate}[label=(\roman*)]
\setlength{\itemsep}{0pt}
\item If $A$ is diagonally dominant, then $A$ is positive semidefinite.
\item If $A$ is diagonally dominant and invertible, then $A$ is positive definite.
\item If $A$ is strictly diagonally dominant, then $A$ is positive definite.
\end{enumerate}
\begin{proof}
Hermitian matrices have real spectrum.
\begin{enumerate}[label=(\roman*)]
\setlength{\itemsep}{0pt}
\item Corollary \ref{cor:strict-diagonal-dominance-invertible} gives $\operatorname{spec}(A)\subseteq[0,\infty)$; hence $A\succeq0$ by Theorem \ref{thm:positive-semidefinite-matrix-equivalences}.
\item If $A$ is invertible, then $0\notin\operatorname{spec}(A)$, so every eigenvalue is positive; hence $A\succ0$.
\item Strict diagonal dominance implies invertibility by Corollary \ref{cor:strict-diagonal-dominance-invertible}; apply (ii).
\end{enumerate}
\hlt{Structural Commentary}: Gersgorin traps the real spectrum in $[0,\infty)$; invertibility removes $0$.\\
\hlt{Key Insight}: For Hermitian matrices, positivity is spectral location.
\end{proof}
\end{corollary}

\begin{theorem} \label{thm:weak-diagonal-dominance-invertible-with-one-strict-row} \hlt{Weak Diagonal Dominance Criterion}\\
Let $A=[a_{ij}]\in\operatorname{Mat}(n,n,\mathbb{C})$ be diagonally dominant with no zero entries.\\
If $\abs{a_{kk}}>R_k^{\mathrm{off}}(A)$ for some $k$, then $A$ is invertible.
\begin{proof}
Assume $A\mathbf{x}=0$ with $\mathbf{x}\neq\mathbf{0}$, and choose $k$ with $\abs{x_k}=\norm{\mathbf{x}}_{\infty}$.\\
Row $k$ gives $\abs{a_{kk}}\norm{\mathbf{x}}_{\infty} =\abs{-\sum\limits_{j\neq k}a_{kj}x_j} \leq\sum\limits_{j\neq k}\abs{a_{kj}}\abs{x_j} \leq R_k^{\mathrm{off}}(A)\norm{\mathbf{x}}_{\infty} \leq \abs{a_{kk}}\norm{\mathbf{x}}_{\infty}$.\\
Hence all inequalities are equalities. Since no $a_{kj}$ vanishes, $\abs{x_j}=\norm{\mathbf{x}}_{\infty}$ for every $j$. Applying the same chain to any row $i$ gives $\abs{a_{ii}}=R'_i(A)$ for every $i$, contradicting the strict row. Thus $\ker(A)=\{\mathbf{0}\}$, so $A$ is invertible.\\
\hlt{Structural Commentary}: A zero vector in the kernel would force equality in every Gersgorin row.\\
\hlt{Key Insight}: One strict row breaks that equality pattern.
\end{proof}
\end{theorem}

\begin{definition} \label{def:increasing-hermitian-eigenvalue-ordering} \hlt{Increasing Hermitian Eigenvalue Order}\\
For a Hermitian matrix $A\in\operatorname{Mat}(n,n,\mathbb{C})$, write $\lambda_1(A)\leq\lambda_2(A)\leq\cdots\leq\lambda_n(A)$ for its eigenvalues counted with algebraic multiplicity. If only one Hermitian matrix is under discussion, we write $\lambda_i$ instead of $\lambda_i(A)$.
\end{definition}

\begin{theorem} \label{thm:rayleigh-quotient-spectral-interval} \hlt{Rayleigh Quotient on a Spectral Interval}\\
Let $A\in\operatorname{Mat}(n,n,\mathbb{C})$ be Hermitian with increasing eigenvalues $\lambda_1\leq\cdots\leq\lambda_n$ and corresponding orthonormal eigenvectors $\mathbf{u}_1,\ldots,\mathbf{u}_n$. For a unit vector, write $\operatorname{Ray}_A(\mathbf{x})\defeq\mathbf{x}^*A\mathbf{x}$. If $1\leq p\leq q\leq n$ and $\mathbf{x}\in\operatorname{span}(\mathbf{u}_p,\ldots,\mathbf{u}_q)$ is a unit vector, then $\lambda_p\leq \operatorname{Ray}_A(\mathbf{x})\leq \lambda_q$.
\begin{proof}
Write $\mathbf{x}=\sum_{i=p}^{q}c_i\mathbf{u}_i$. Orthonormality gives $\sum_{i=p}^{q}\abs{c_i}^2=1$, and $\mathbf{x}^*A\mathbf{x} =\sum\limits_{i=p}^{q}\lambda_i\abs{c_i}^2$.\\
This is a convex combination of $\lambda_p,\ldots,\lambda_q$, so it lies between the smallest and largest of them.\\
\hlt{Structural Commentary}: In an eigenbasis, the Rayleigh quotient is an eigenvalue average.\\
\hlt{Key Insight}: Spectral coordinates turn a quadratic form into weighted scalar data.
\end{proof}
\end{theorem}

\begin{corollary} \label{cor:rayleigh-extrema-and-diagonal-entry-bounds} \hlt{Rayleigh Extrema and Diagonal Entry Bounds}\\
Let $A\in\operatorname{Mat}(n,n,\mathbb{C})$ be Hermitian.
\begin{enumerate}[label=(\roman*)]
\setlength{\itemsep}{0pt}
\item If $\mathcal{U}_{p,q}=\operatorname{span}(\mathbf{u}_p,\ldots,\mathbf{u}_q)$, then $\min\limits_{\substack{\mathbf{x}\in\mathcal{U}_{p,q}\\ \norm{\mathbf{x}}_2=1}}\mathbf{x}^*A\mathbf{x}=\lambda_p(A)$, and $\max\limits_{\substack{\mathbf{x}\in\mathcal{U}_{p,q}\\ \norm{\mathbf{x}}_2=1}}\mathbf{x}^*A\mathbf{x}=\lambda_q(A).$
\item For every unit vector $\mathbf{x}$, $\lambda_1(A)\leq \mathbf{x}^*A\mathbf{x}\leq\lambda_n(A)$.
\item If $A=[a_{ij}]$, then $\lambda_1(A)\leq a_{ii}\leq\lambda_n(A)$ for $i=1,\ldots,n$.
\end{enumerate}
\begin{proof}
\begin{enumerate}[label=(\roman*)]
\setlength{\itemsep}{0pt}
\item Theorem \ref{thm:rayleigh-quotient-spectral-interval} gives the inequalities on $\mathcal{U}_{p,q}$, with equality at $\mathbf{u}_p$ and $\mathbf{u}_q$.
\item This is (i) with $p=1$ and $q=n$.
\item Apply (ii) to $\mathbf{x}=\mathbf{e}_i$, since $\mathbf{e}_i^*A\mathbf{e}_i=a_{ii}$.
\end{enumerate}
\hlt{Structural Commentary}: Rayleigh extrema on extreme eigendirections; diagonal entries are coordinate quotients.\\
\hlt{Key Insight}: The whole numerical range of a Hermitian matrix sits inside its spectral interval.
\end{proof}
\end{corollary}

\begin{figure}[H]
  \centering
  \begin{tikzpicture}[
  >=Stealth,
  scale=0.88,
  axis/.style={->, gray!65, line width=0.65pt},
  eig/.style={circle, fill=blue!65!black, inner sep=1.7pt},
  rq/.style={circle, fill=red!75!black, inner sep=2.2pt},
  lab/.style={font=\small, align=center, fill=white, inner sep=1pt},
  tiny/.style={font=\small, align=center}
]
  \draw[axis] (-0.4,0) -- (7.4,0);
  \foreach \x/\labx in {0.8/{\lambda_1},2.2/{\lambda_p},5.0/{\lambda_q},6.4/{\lambda_n}} {
    \draw[gray!55] (\x,0.09) -- (\x,-0.09);
    \node[eig] at (\x,0) {};
    \node[tiny, below=4pt] at (\x,0) {$\labx$};
  }
  \draw[green!55!black, line width=1.25pt] (2.2,0.55) -- (5.0,0.55);
  \node[lab, above=4pt] at (3.6,0.55) {$\mathcal{U}_{p,q}$};
  \node[rq] at (3.7,0) {};
  \node[lab, below=16pt] at (3.7,0) {$\operatorname{Ray}_A(\mathbf{x})$};
\end{tikzpicture}

  \caption{For $\mathbf{x}\in\mathcal{U}_{p,q}$, $\operatorname{Ray}_A(\mathbf{x})$ is a convex combination of $\lambda_p,\ldots,\lambda_q$.}
  \label{fig:rayleigh-quotient-convex-combination}
\end{figure}

\begin{theorem} \label{thm:courant-fischer-rayleigh-minmax} \hlt{Rayleigh Minimax Principle}\\
Let $A=A^*\in\operatorname{Mat}(n,n,\mathbb{C})$ have eigenvalues $\lambda_1\leq\cdots\leq\lambda_n$, and set $R_A(\mathbf{x})=\mathbf{x}^*A\mathbf{x}/(\mathbf{x}^*\mathbf{x})$ for $\mathbf{x}\neq\mathbf{0}$.\\
For $1\leq k\leq n$, $\lambda_k =\min\limits_{\substack{\mathcal{S}\subseteq\mathbb{C}^n\\ \dim\mathcal{S}=k}}\; \max\limits_{\mathbf{0}\neq\mathbf{x}\in\mathcal{S}} R_A(\mathbf{x}) =\max\limits_{\substack{\mathcal{S}\subseteq\mathbb{C}^n\\ \dim\mathcal{S}=n-k+1}}\; \min\limits_{\mathbf{0}\neq\mathbf{x}\in\mathcal{S}} R_A(\mathbf{x})$.
\begin{proof}
Let $\mathbf{u}_1,\ldots,\mathbf{u}_n$ be an orthonormal eigenbasis in increasing order.\\
Set $E_k=\operatorname{span}(\mathbf{u}_1,\ldots,\mathbf{u}_k)$; $F_k=\operatorname{span}(\mathbf{u}_k,\ldots,\mathbf{u}_n)$.\\
By Theorem \ref{thm:rayleigh-quotient-spectral-interval}, $\max_{\mathbf{0}\neq\mathbf{x}\in E_k}R_A(\mathbf{x})=\lambda_k$ and $\min_{\mathbf{0}\neq\mathbf{x}\in F_k}R_A(\mathbf{x})=\lambda_k$.\\
If $\dim\mathcal{S}=k$, then $\mathcal{S}\cap F_k\neq\{\mathbf{0}\}$, so $\max_{\mathbf{0}\neq\mathbf{x}\in\mathcal{S}}R_A(\mathbf{x})\geq\lambda_k$; equality is attained by $E_k$.\\
If $\dim\mathcal{T}=n-k+1$, then $\mathcal{T}\cap E_k\neq\{\mathbf{0}\}$, so $\min_{\mathbf{0}\neq\mathbf{x}\in\mathcal{T}}R_A(\mathbf{x})\leq\lambda_k$; equality is attained by $F_k$.\\
\hlt{Structural Commentary}: Dimension forces intersections with spectral flags.\\
\hlt{Key Insight}: Minimax is Rayleigh's quotient plus dimension counting.
\end{proof}
\end{theorem}

\begin{method} \label{meth:rayleigh-eigenbasis-method} \hlt{Rayleigh Eigenbasis Method}\\
To bound a Hermitian quadratic form:
\begin{enumerate}[label=(\roman*)]
\setlength{\itemsep}{0pt}
\item diagonalise $A$ with an orthonormal eigenbasis;
\item expand the unit vector in that basis;
\item rewrite $\mathbf{x}^*A\mathbf{x}$ as $\sum_i\lambda_i\abs{c_i}^2$;
\item use the eigenvalue range allowed by the support of $\mathbf{x}$.
\end{enumerate}
\end{method}

\begin{theorem} \label{thm:rank-one-psd-eigenvalue-interlacing} \hlt{Rank-One Positive Perturbation Interlacing}\\
Let $A,E\in\operatorname{Mat}(n,n,\mathbb{C})$ be Hermitian, with $E\succeq0$ and $\operatorname{rank}(E)\leq1$.\\
Then $\lambda_i(A)\leq\lambda_i(A+E)\leq\lambda_{i+1}(A)\quad(1\leq i<n)$, $\lambda_n(A)\leq\lambda_n(A+E)$\\
Equivalently, $\lambda_i(A-E)\leq\lambda_i(A)\leq\lambda_{i+1}(A-E)\quad(1\leq i<n)$, $\lambda_n(A-E)\leq\lambda_n(A)$.
\begin{proof}
Let $\mathbf{u}_j$ and $\mathbf{v}_j$ be orthonormal eigenvectors of $A$ and $A+E$, respectively, in increasing eigenvalue order.\\
For the lower bound, set $\mathcal{U}=\operatorname{span}(\mathbf{u}_i,\ldots,\mathbf{u}_n)$, $\mathcal{V}=\operatorname{span}(\mathbf{v}_1,\ldots,\mathbf{v}_i)$. Since $\dim\mathcal{U}+\dim\mathcal{V}=n+1$, choose a unit $\mathbf{x}\in\mathcal{U}\cap\mathcal{V}$. Then $\lambda_i(A)\leq \mathbf{x}^*A\mathbf{x}\leq \mathbf{x}^*(A+E)\mathbf{x}\leq\lambda_i(A+E)$ by Theorem \ref{thm:rayleigh-quotient-spectral-interval} and $E\succeq0$.\\
For the upper bound, set $\mathcal{U}=\operatorname{span}(\mathbf{u}_1,\ldots,\mathbf{u}_{i+1})$, $\mathcal{V}=\operatorname{span}(\mathbf{v}_i,\ldots,\mathbf{v}_n)$, $\mathcal{W}=\ker(E)$. Here $\dim(\mathcal{U}\cap\mathcal{V})\geq2$ and $\dim\mathcal{W}\geq n-1$, so choose a unit $\mathbf{x}\in\mathcal{U}\cap\mathcal{V}\cap\mathcal{W}$. Thus $\lambda_i(A+E)\leq \mathbf{x}^*(A+E)\mathbf{x} =\mathbf{x}^*A\mathbf{x} \leq \lambda_{i+1}(A)$.\\
The statements for $A-E$ follow by applying the proved result to $A-E$ and $E$.\\
\hlt{Structural Commentary}: Interlacing comes from test vectors in overlapping spectral subspaces.\\
\hlt{Key Insight}: Rank one leaves a large kernel, so the upper-bound vector can avoid the perturbation.
\end{proof}
\end{theorem}

\begin{corollary} \label{cor:rank-two-one-positive-one-negative-interlacing} \hlt{Rank-Two Indefinite Interlacing}\\
Let $A,H\in\operatorname{Mat}(n,n,\mathbb{C})$ be Hermitian, $n\geq2$. Suppose $\operatorname{rank}(H)\leq2$ and $H$ has at most one positive and one negative eigenvalue. Then
\begin{align}
\lambda_1(A+H)\leq\lambda_2(A),\qquad
\lambda_{i-1}(A)\leq\lambda_i(A+H)\leq\lambda_{i+1}(A)\quad(2\leq i<n),
\qquad
\lambda_{n-1}(A)\leq\lambda_n(A+H).
\nonumber
\end{align}
\begin{proof}
Write $H=E_+-E_-$ with $E_\pm\succeq0$ and $\operatorname{rank}(E_\pm)\leq1$. Set $B=A-E_-$; then $A=B+E_-$ and $A+H=B+E_+$. For $2\leq i<n$, rank-one interlacing gives $\lambda_{i-1}(A)\leq\lambda_i(B)\leq\lambda_i(A+H)\leq\lambda_{i+1}(B)\leq\lambda_{i+1}(A)$.\\
The endpoints follow similarly: $\lambda_1(A+H)\leq\lambda_2(B)\leq\lambda_2(A)$, and $\lambda_{n-1}(A)\leq\lambda_n(B)\leq\lambda_n(A+H)$.\\
\hlt{Structural Commentary}: Such an $H$ is a difference of two rank-one PSD perturbations.\\
\hlt{Key Insight}: Interlace once with $E_-$, then once with $E_+$.
\end{proof}
\end{corollary}

\begin{figure}[H]
  \centering
  \begin{tikzpicture}[
  >=Stealth,
  scale=1,
  axis/.style={->, gray!65, line width=0.65pt},
  adot/.style={circle, fill=blue!65!black, inner sep=1.8pt},
  bdot/.style={circle, fill=red!70!black, inner sep=1.8pt},
  cdot/.style={circle, fill=green!45!black, inner sep=1.8pt},
  guide/.style={dashed, gray!55, line width=0.55pt},
  lab/.style={font=\small, align=center, fill=white, inner sep=1pt},
  tiny/.style={font=\small, align=center, fill=white, inner sep=1pt}
]
  \draw[axis] (0,1.05) -- (3.55,1.05) node[right, lab] {$A+E$};
  \draw[axis] (0,0) -- (3.55,0) node[right, lab] {$A$};

  \foreach \x in {0.85,2.55} {
    \node[adot] at (\x,0) {};
    \draw[guide] (\x,1.22) -- (\x,-0.24);
  }
  \node[bdot] at (1.45,1.05) {};
  \draw[<->, blue!55!black] (0.85,0.34) -- (2.55,0.34);
  \node[tiny, below=5pt] at (0.85,0) {$\lambda_i(A)$};
  \node[tiny, below=5pt] at (2.55,0) {$\lambda_{i+1}(A)$};
  \node[tiny, above=5pt] at (1.45,1.05) {$\lambda_i(A+E)$};

  \draw[axis] (5.1,1.05) -- (8.95,1.05) node[right, lab] {$B$};
  \draw[axis] (5.1,0) -- (8.95,0) node[right, lab] {$A$};

  \foreach \x in {6.0,7.9} {
    \node[adot] at (\x,0) {};
    \draw[guide] (\x,1.22) -- (\x,-0.24);
  }
  \node[cdot] at (6.65,1.05) {};
  \draw[<->, green!45!black] (6.0,0.34) -- (7.9,0.34);
  \node[tiny, below=5pt] at (6.0,0) {$\lambda_i(A)$};
  \node[tiny, below=5pt] at (7.9,0) {$\lambda_{i+m}(A)$};
  \node[tiny, above=5pt] at (6.65,1.05) {$\lambda_i(B)$};
\end{tikzpicture}

  \caption{Rank-one and principal-submatrix interlacing place new eigenvalues between old ones.}
  \label{fig:eigenvalue-interlacing-ladder}
\end{figure}

\begin{theorem} \label{thm:cauchy-principal-submatrix-interlacing} \hlt{Cauchy Interlacing: One Principal Deletion}\\
Let $A=
\begin{bmatrix}
B & \mathbf{y}\\
\mathbf{y}^* & c
\end{bmatrix}
\in\operatorname{Mat}(n+1,n+1,\mathbb{C})$ be Hermitian, with $B\in\operatorname{Mat}(n,n,\mathbb{C})$.\\
Then, for $1\leq i\leq n$, $\lambda_i(A)\leq\lambda_i(B)\leq\lambda_{i+1}(A)$.
\begin{proof}
Let $\mathbf{u}_j$ and $\mathbf{v}_j$ be orthonormal eigenvectors of $A$ and $B$, respectively, in increasing eigenvalue order. Identify $\mathbb{C}^n$ with $\{(\mathbf{z},0)\}\subseteq\mathbb{C}^{n+1}$, so $(\mathbf{z},0)^*A(\mathbf{z},0)=\mathbf{z}^*B\mathbf{z}$.\\
For the lower bound, set $\mathcal{S}=\operatorname{span}((\mathbf{v}_1,0),\ldots,(\mathbf{v}_i,0))$ and $\mathcal{T}=\operatorname{span}(\mathbf{u}_i,\ldots,\mathbf{u}_{n+1})$. Their dimensions are $i$ and $n+2-i$, so $\dim\mathcal{S}+\dim\mathcal{T}>n+1$ and $\mathcal{S}\cap\mathcal{T}$ contains a nonzero vector.\\
Choose a unit $\mathbf{x}=(\mathbf{z},0)$ in the intersection. The spectral expansion in the $\mathbf{u}_j$ gives $\mathbf{x}^*A\mathbf{x}\geq\lambda_i(A)$, while the expansion in the $\mathbf{v}_j$ gives $\mathbf{z}^*B\mathbf{z}\leq\lambda_i(B)$. Hence $\lambda_i(A)\leq\lambda_i(B)$.\\
For the upper bound, take $\mathcal{S}'=\operatorname{span}((\mathbf{v}_i,0),\ldots,(\mathbf{v}_n,0))$ and $\mathcal{T}'=\operatorname{span}(\mathbf{u}_1,\ldots,\mathbf{u}_{i+1})$. Again the dimensions sum to $n+2$, so choose a unit $\mathbf{x}=(\mathbf{z},0)\in\mathcal{S}'\cap\mathcal{T}'$.\\
Now $\mathbf{z}^*B\mathbf{z}\geq\lambda_i(B)$ and $\mathbf{x}^*A\mathbf{x}\leq\lambda_{i+1}(A)$, giving $\lambda_i(B)\leq\lambda_{i+1}(A)$.\\
\hlt{Structural Commentary}: A principal submatrix is a coordinate slice of the larger quadratic form.\\
\hlt{Key Insight}: Deleting one coordinate costs one dimension, so eigenvalues shift by at most one index.
\end{proof}
\end{theorem}

\begin{corollary} \label{cor:singular-values-interlace-after-deleting-row-or-column} \hlt{Singular Values Interlace After Deleting a Row or Column}\\
Let $B$ be obtained from $A\in\operatorname{Mat}(n,n,\mathbb{C})$ by deleting one row or column.\\
With singular values ordered decreasingly, $\sigma_i(A)\geq\sigma_i(B)\geq\sigma_{i+1}(A)\quad(1\leq i<n)$.
\begin{proof}
For a row deletion, permute rows so that $PA=\begin{bmatrix}B\\ \mathbf{y}^*\end{bmatrix}$. \\
Then $PA$ has the same singular values as $A$, and $BB^*$ is a principal submatrix of $(PA)(PA)^*$.\\
Cauchy interlacing applied to these Gram matrices gives the result after taking square roots and reversing the increasing eigenvalue order. A column deletion follows by applying the row case to $A^*$.\\
\hlt{Structural Commentary}: Row deletion makes $BB^*$ a principal submatrix of a permuted $AA^*$.\\
\hlt{Key Insight}: Singular-value interlacing is Cauchy interlacing applied to Gram matrices.
\end{proof}
\end{corollary}

\begin{corollary} \label{cor:cauchy-interlacing-general-principal-submatrix} \hlt{Cauchy Interlacing for Principal Submatrices}\\
Let $B\in\operatorname{Mat}(n,n,\mathbb{C})$ be a principal submatrix of Hermitian $A\in\operatorname{Mat}(n+m,n+m,\mathbb{C})$.\\
Then, for $1\leq i\leq n$, $\lambda_i(A)\leq\lambda_i(B)\leq\lambda_{i+m}(A)$.
\begin{proof}
After a permutation similarity, $B$ is leading.\\
Delete the other $m$ coordinates one at a time and apply Theorem \ref{thm:cauchy-principal-submatrix-interlacing} at each step.\\
\hlt{Structural Commentary}: Principal submatrices come from repeated coordinate deletion.\\
\hlt{Key Insight}: Each deletion shifts eigenvalue indices by at most one.
\end{proof}
\end{corollary}

\begin{theorem} \label{thm:schur-majorization-eigenvalues-by-diagonal} \hlt{Schur Majorization of the Diagonal}\\
Let $A\in\operatorname{Mat}(n,n,\mathbb{C})$ be Hermitian, with increasing eigenvalues $\lambda_i(A)$ and increasing diagonal entries $d_i$.\\
Then $\sum\limits_{i=1}^{k}\lambda_i(A)\leq \sum\limits_{i=1}^{k}d_i$, $(1\leq k\leq n)$ with equality for $k=n$.
\begin{proof}
Permute coordinates so that $a_{ii}=d_i$, and let $A_k$ be the leading $k\times k$ principal submatrix.\\
Cauchy interlacing gives $\lambda_i(A)\leq\lambda_i(A_k)$ for $1\leq i\leq k$, hence $\sum\limits_{i=1}^{k}\lambda_i(A)\leq\sum\limits_{i=1}^{k}\lambda_i(A_k)=\operatorname{tr}(A_k)=\sum\limits_{i=1}^{k}d_i$.\\
For $k=n$, this is $\operatorname{tr}(A)=\operatorname{tr}(A)$.\\
\hlt{Structural Commentary}: Ordered diagonal sums are traces of leading principal blocks.\\
\hlt{Key Insight}: Diagonal majorization is Cauchy interlacing plus trace.
\end{proof}
\end{theorem}

\begin{method} \label{meth:eigenvalue-interlacing-method} \hlt{Eigenvalue Interlacing Method}\\
To prove or use an interlacing estimate:
\begin{enumerate}[label=(\roman*)]
\setlength{\itemsep}{0pt}
\item order Hermitian eigenvalues increasingly;
\item identify dimension lost: rank-one perturbation, one deleted coordinate, $m$-coordinate principal submatrix;
\item choose test vectors by intersecting spectral subspaces with the unaffected subspace;
\item convert the Rayleigh quotient bounds into eigenvalue inequalities.
\end{enumerate}
\end{method}

\begin{theorem} \label{thm:weyl-eigenvalue-inequalities-hermitian-sums} \hlt{Weyl Eigenvalue Inequalities}\\
Let $A,B\in\operatorname{Mat}(n,n,\mathbb{C})$ be Hermitian. For $1\leq i\leq n$,
\begin{align*}
\lambda_i(A+B) &\leq \lambda_{i+j}(A)+\lambda_{n-j}(B),
& 0 &\leq j\leq n-i,\\
\lambda_{i-j+1}(A)+\lambda_j(B) &\leq \lambda_i(A+B),
& 1 &\leq j\leq i.
\end{align*}
\begin{proof}
Let $\mathbf{u}_k,\mathbf{v}_k,\mathbf{w}_k$ be increasing eigenbases for $A,A+B,B$.\\
Fix $i$ and $0\leq j\leq n-i$, and set $\mathcal{U}=\operatorname{span}(\mathbf{u}_1,\ldots,\mathbf{u}_{i+j})$, $\mathcal{V}=\operatorname{span}(\mathbf{v}_i,\ldots,\mathbf{v}_n)$, $\mathcal{W}=\operatorname{span}(\mathbf{w}_1,\ldots,\mathbf{w}_{n-j})$.\\
Their dimensions sum to $2n+1$, so choose a unit $\mathbf{x}\in\mathcal{U}\cap\mathcal{V}\cap\mathcal{W}$.\\
Then $\lambda_i(A+B)\leq \mathbf{x}^*(A+B)\mathbf{x} =\mathbf{x}^*A\mathbf{x}+\mathbf{x}^*B\mathbf{x} \leq \lambda_{i+j}(A)+\lambda_{n-j}(B)$.\\
The lower bound follows by applying this upper bound to $-A$ and $-B$ and using $\lambda_k(-C)=-\lambda_{n-k+1}(C)$.\\
\hlt{Structural Commentary}: Weyl bounds come from intersecting the spectral flags of $A$, $B$, and $A+B$.\\
\hlt{Key Insight}: One shared test vector controls all three Rayleigh quotients.
\end{proof}
\end{theorem}

\begin{method} \label{meth:weyl-inequality-method} \hlt{Weyl Bound Method}\\
To bound eigenvalues of $A+B$:
\begin{enumerate}[label=(\roman*)]
\setlength{\itemsep}{0pt}
\item choose the target index $i$;
\item use the upper Weyl family with $j=0,\ldots,n-i$ to get upper bounds;
\item use the lower Weyl family with $j=1,\ldots,i$ to get lower bounds;
\item for positive semidefinite or low-rank $B$, insert the known zero or nonnegative eigenvalues to recover monotonicity or interlacing.
\end{enumerate}
\end{method}

\begin{theorem} \label{thm:brauer-rank-one-eigenvalue-shift} \hlt{Brauer Rank-One Eigenvalue Shift}\\
Let $A\in\operatorname{Mat}(n,n,\mathbb{C})$, let $A\mathbf{x}=\lambda\mathbf{x}$ with $\mathbf{x}\neq\mathbf{0}$, and let $\mathbf{y}\in\mathbb{C}^n$. If $\lambda,\lambda_2,\ldots,\lambda_n$ are the eigenvalues of $A$ with multiplicity, then the eigenvalues of $A+\mathbf{x}\mathbf{y}^*$ are $\lambda+\mathbf{y}^*\mathbf{x},\lambda_2,\ldots,\lambda_n$.\\
Moreover, $\mathbf{x}$ is an eigenvector of $A+\mathbf{x}\mathbf{y}^*$ corresponding to $\lambda+\mathbf{y}^*\mathbf{x}$.
\begin{proof}
Put $\mathbf{u}=\mathbf{x}/\norm{\mathbf{x}}_2$. By Schur triangularisation, choose a unitary $U=[\mathbf{u}\ U_2]$ such that $U^*AU=\begin{bmatrix}\lambda&*\\0&T'\end{bmatrix}$, where the diagonal entries of $T'$ are $\lambda_2,\ldots,\lambda_n$. Also $U^*(\mathbf{x}\mathbf{y}^*)U
=
\begin{bmatrix}
\mathbf{y}^*\mathbf{x}&*\\
0&0
\end{bmatrix}$.
Therefore $U^*(A+\mathbf{x}\mathbf{y}^*)U$ is upper triangular with diagonal entries $\lambda+\mathbf{y}^*\mathbf{x},\lambda_2,\ldots,\lambda_n$. Finally, $(A+\mathbf{x}\mathbf{y}^*)\mathbf{x}=A\mathbf{x}+\mathbf{x}(\mathbf{y}^*\mathbf{x})=(\lambda+\mathbf{y}^*\mathbf{x})\mathbf{x}$.
\hlt{Structural Commentary}: A rank-one eigenvector update changes only the first Schur diagonal entry.\\
\hlt{Key Insight}: Align the update with the first coordinate; only one diagonal term shifts.
\end{proof}
\end{theorem}

\begin{figure}[H]
  \centering
  \begin{tikzpicture}[
  >=Stealth,
  every node/.style={font=\small},
  mat/.style={draw=gray!55, fill=gray!5, rounded corners=2pt, align=center, inner sep=6pt},
  shiftbox/.style={draw=blue!60!black, fill=blue!7, rounded corners=2pt, align=center, inner sep=6pt},
  arr/.style={->, line width=0.75pt, blue!70!black}
]
  \node[mat] (schur) at (0,0)
    {$U^*AU=\begin{bmatrix}\lambda&*\\0&T'\end{bmatrix}$};
  \node[shiftbox] (rank) at (4,0)
    {$U^*\mathbf{x}\mathbf{y}^*U=\begin{bmatrix}\mathbf{y}^*\mathbf{x}&*\\0&0\end{bmatrix}$};
  \node[mat] (sum) at (8.2,0)
    {$\begin{bmatrix}\lambda+\mathbf{y}^*\mathbf{x}&*\\0&T'\end{bmatrix}$};

  \draw[arr] (schur) -- (rank);
  \draw[arr] (rank) -- (sum);
\end{tikzpicture}

  \caption{Brauer's update shifts one Schur diagonal entry.}
  \label{fig:brauer-rank-one-shift}
\end{figure}

\begin{corollary} \label{cor:damped-rank-one-eigenvalue-reset} \hlt{Damped Rank-One Eigenvalue Reset}\\
Let $A\in\operatorname{Mat}(n,n,\mathbb{C})$ have eigenpair $(\lambda,\mathbf{x})$, let $\mathbf{y}^*\mathbf{x}=1$, and let $\tau\in\mathbb{C}$. If $\lambda,\lambda_2,\ldots,\lambda_n$ are the eigenvalues of $A$, then $A_{\tau}\defeq \tau A+(1-\tau)\lambda\mathbf{x}\mathbf{y}^*$ has eigenvalues $\lambda,\tau\lambda_2,\ldots,\tau\lambda_n$.
\begin{proof}
The eigenvalues of $\tau A$ are $\tau\lambda,\tau\lambda_2,\ldots,\tau\lambda_n$. Since $A_{\tau}=\tau A+\mathbf{x}((1-\overline{\tau})\overline{\lambda}\mathbf{y})^*$, Theorem \ref{thm:brauer-rank-one-eigenvalue-shift} shifts only $\tau\lambda$ by $(1-\tau)\lambda\mathbf{y}^*\mathbf{x}=(1-\tau)\lambda$. Thus $\tau\lambda$ becomes $\lambda$, while the other eigenvalues stay $\tau\lambda_i$.\\
\hlt{Structural Commentary}: Damping scales the spectrum; Brauer's update restores one eigenvalue.\\
\hlt{Key Insight}: Brauer's theorem changes one spectral coordinate without moving the rest.
\end{proof}
\end{corollary}
